\documentclass[11pt]{article}

\usepackage{amsmath,amssymb}
\usepackage{xurl}
\usepackage[hidelinks]{hyperref}
\setlength{\emergencystretch}{2em}

% Shared commands for notes in docs/paper-gaps/.
%
% Two halves.  The link commands below are project identity: OWNER/REPO is the
% placeholder that scripts/bootstrap_project.py replaces with this project's
% GitHub slug and Pages site.  Everything after them is NOTATION, and notation
% belongs to the source paper: the definitions here are an example set, and a
% project replaces them with the macros of its own paper's style file.  The one
% rule that never changes: a command defined here must mean exactly what the
% source paper's macro means, or a note silently says something else.

\newcommand{\Lean}{\textsc{Lean}}
\newcommand{\leanid}[1]{\textsf{#1}}
\newcommand{\ghissue}[1]{\href{https://github.com/OWNER/REPO/issues/#1}{GitHub issue \##1}}
\newcommand{\ghpr}[1]{\href{https://github.com/OWNER/REPO/pull/#1}{PR \##1}}
% Local issue tracker (issues/NNNN-*.md in this repository); no remote URL.
\newcommand{\localissue}[1]{issue \##1 (\path{issues/})}
\newcommand{\doclink}[1]{\href{https://OWNER.github.io/REPO/docs/#1.html}{\path{#1}}}

% --- Notation: replace with the source paper's own macros --------------------
% \F, \N, \C, \R, \tr, \ot, \eps, \E, \bx, \bu, \bv, \by

\newcommand{\F}{\mathbb{F}}
\newcommand{\N}{\mathbb{N}}
\newcommand{\C}{\mathbb{C}}
\newcommand{\R}{\mathbb{R}}
\newcommand{\tr}{\mathrm{tr}}
\newcommand{\eps}{\epsilon}
\newcommand{\ot}{\otimes}
\newcommand{\E}{\mathop{\bf E\/}}

% bold random variables
\newcommand{\bu}{\boldsymbol{u}}
\newcommand{\bv}{\boldsymbol{v}}
\newcommand{\bx}{{\boldsymbol{x}}}
\newcommand{\by}{\boldsymbol{y}}

% T1 so literal < and > in \texttt placeholders typeset as themselves
\usepackage[T1]{fontenc}

% bra-ket
\usepackage{braket}

% --- Example structured spaces ---
\newcommand{\polyfunc}[3]{\mathcal{P}(#1, #2, #3)}
\newcommand{\polymeas}[3]{\mathrm{PolyMeas}(#1, #2, #3)}
\newcommand{\polysub}[3]{\mathrm{PolySub}(#1, #2, #3)}

% Approximation relations (paper uses \approx_\eps and \simeq_\zeta)
\newcommand{\approxeps}{\approx_{\varepsilon}}
\newcommand{\simgeq}{\simeq_{\zeta}}


\newcommand{\dbin}{d^{\mathrm{bin}}}
\newcommand{\Id}{\mathrm{Id}}
\providecommand{\gapnote}[2]{\par\noindent\textbf{Verdict:} #1 (#2).\par}
\renewcommand{\ghissue}[1]{%
  \href{https://github.com/Dengnifer/MIPStarRE-QPBT/issues/#1}{GitHub issue \##1}}

\title{Distance Normalization in the Quantum Linearity Theorem}
\author{}
\date{2026-09-23}

\begin{document}
\maketitle
\gapnote{false-source}{resolved}
\noindent Resolved by documentation; register status
\texttt{documented-deviation}. The printed operator bound remains refuted.

\section{At a glance}

\begin{itemize}
  \item \textbf{Difficulty.} Medium.  The factor-of-two calculation is
  elementary, but the counterexample must remain valid after the ancillary
  extension allowed by the theorem.
  \item \textbf{Estimated weight.} The normalization identities and the
  Fourier--Naimark theorem were already proved. The additional verification
  uses roughly 250 lines for the printed propositions and the four-point
  obstruction, including arbitrary ancillary dimensions, and 35 lines for
  the corrected theorem without numerical domain restrictions.
  \item \textbf{Mathlib/project split.} About 60\% is direct finite-matrix and
  finite-sum algebra; the remaining 40\% is the project-local Boolean Fourier
  and state-dependent distance interface.
  \item \textbf{Key Mathlib inputs.} Finite sums on the Boolean cube, complex
  matrices and conjugate transpose, cyclicity of trace, and simultaneous
  spectral decomposition for commuting Hermitian matrices.
\end{itemize}
The squared operator bound $2\delta$ is valid, and $2$ is the least universal
multiplicative constant, even for positive index length and $0<\delta<1$.
The sharpness argument below uses a family of Boolean quadratic functions;
the four-point example alone only forces $4/3$. This is not a claim that the
Fourier--Naimark construction minimizes the error for each input. This note
closes the normalization discrepancy as a documented constant/convention
deviation; it does not certify the refuted printed operator bound. The
separate state-preserving padding question remains open. The corrected
existence theorem, exact distance/defect identity, and four-point optimum
are already formalized; the asymptotic sharpness argument is mathematical
and is not claimed as a \Lean{} theorem.%
\footnote{The normalization evidence of issue \#694 was merged in
\href{https://github.com/Dengnifer/MIPStarRE-QPBT/pull/699}{PR~\#699},
commit \texttt{317aa9f4eb625f502043f57398e43f8e80201bf6}, before this
documentation update. Under the owner's 2026-09-22 instruction, issue
\href{https://github.com/Dengnifer/MIPStarRE-QPBT/issues/718}{\#718}
closes only this row by documentation, superseding its earlier strict-adoption
completion requirement without changing any theorem or asserting the printed
claim. The earlier audit and budget evidence remain in
\path{audits/2026-09-22_linearity-normalization.md}.}

\section{Key theorem forms}

Let $S$ and $T$ be binary observables and let $\rho$ be a positive
semidefinite operator of trace one.  We distinguish the state-dependent
operator distance from the distance between their associated binary projective
measurements:
\begin{align*}
  d_\rho(S,T)^2
    &:= \Re\tr\big((S-T)^\dagger(S-T)\rho\big), \\
  \dbin_\rho(S,T)^2
    &:= \sum_{s\in\{+,-\}}
      \Re\tr\big((P_S^s-P_T^s)^\dagger
                       (P_S^s-P_T^s)\rho\big),
  \qquad P_S^\pm := \frac{\Id\pm S}{2}.
\end{align*}
The normalization and correlation identities are
\begin{align}
  \dbin_\rho(S,T)^2 &= \frac12 d_\rho(S,T)^2,
    \tag{normalization}\label{eq:normalization}\\
  d_\rho(S,T)^2
    &= 2-2\Re\tr(ST\rho),
    \tag{operator correlation}\label{eq:operator-correlation}\\
  \dbin_\rho(S,T)^2
    &= 1-\Re\tr(ST\rho).
    \tag{binary correlation}\label{eq:binary-correlation}
\end{align}
Consequently, the conclusion delivered by the quantum linearity argument can
be written either as
\[
  \mathbb E_a\dbin_{\rho'}(\mathcal A(a),A(a))^2\leq\delta
\]
or, equivalently, as
\[
  \mathbb E_a d_{\rho'}(\mathcal A(a),A(a))^2\leq2\delta.
\]
These are not equivalent to the same bound $\delta$ for the mean squared
state-dependent operator distance.%
\footnote{The formal declarations are \leanid{binaryObservableEffect},
\leanid{binaryObservableDistSq},
\leanid{binaryObservableDistSq\_eq\_stateDepDistSq\_div\_two},
\leanid{stateDepDistSq\_eq\_two\_sub\_two\_mul\_correlation}, and
\leanid{binaryObservableDistSq\_eq\_one\_sub\_correlation} in
\doclink{MIPStarRE/QPBT/Combining/Linearity/BooleanFourier.lean}.}

More precisely, fix a nonnegative integer $t$, a real number $\delta$, a
finite-dimensional complex Hilbert space $\mathcal H$, a density operator
$\rho$ on it, and Hermitian involutions $A(a)$ for $a\in\mathbb F_2^t$.
Write
\[
 C=\mathbb E_{a,b}\Re\tr(A(a)A(b)A(a+b)\rho).
\]
The printed theorem asserts that $C\geq1-\delta$ implies the existence of
a finite-dimensional space $\mathcal H'$, a unit vector $\mathrm{anc}$ in it,
and Hermitian involutions $\mathcal A(a)$ on $\mathcal H\otimes\mathcal H'$
such that $\mathcal A(a)\mathcal A(b)=\mathcal A(a+b)$ for every pair and
$\mathbb E_a d_{\rho'}(\mathcal A(a),A(a)\otimes\Id)^2\leq\delta$.
Here $\rho'=\rho\otimes|\mathrm{anc}\rangle\langle\mathrm{anc}|$; the
extension is chosen after the input family. The correction changes only
the final bound to $2\delta$. No commutation, upper error bound, or prescribed
ancilla is assumed. Taking real parts makes explicit a real quantity:
conjugation reverses the triple product, and exchanging $a$ with $a+b$
preserves the uniform average. The source's ambient spaces are finite
dimensional by its Section~3.1, line~732.

\section{The source calculation}

Natarajan and Vidick define the distance between two POVMs by summing the
squared state-dependent distances between corresponding square roots
\cite[Section~3.1]{Natarajan2017Linearity}.  For projective measurements the
square roots are the effects themselves.  In the binary case their calculation
expands $P_S^\pm=(\Id\pm S)/2$ and obtains
\[
  \dbin_\rho(S,T)^2
    = 1-\frac12\tr_\rho(ST+TS)
    = \frac12\tr_\rho((S-T)^2).
  \tag{source equation (3)}
\]
This is equation~(3) at lines~903--910 of the in-repository source mirror.%
\footnote{See
\path{references/nv-paper/fullpaper.tex}, lines~866--912.  The quantum
linearity theorem and its Fourier--Naimark proof are at lines~1074--1112.}
The next line identifies the final expression with the squared state-dependent
operator distance.  The operator-distance definition at lines~873--875 is
printed with $\rho$ occurring twice,
\[
  d_\rho(S,T)^2 = \tr_\rho\big((S-T)^\dagger(S-T)\rho\big),
  \qquad \tr_\rho(X) := \tr(\rho X),
\]
and we read the duplication as typographical, keeping the single $\rho$ of the
definition of $d$ above.  Under that reading the squared state-dependent
operator distance is $\tr_\rho((S-T)^2)$ rather than one half of that
quantity.  Thus the calculation through line~910 is correct and the equality
on line~911 loses a factor of two.

The correction is also sufficient if the duplicated density is retained
literally. A density operator has eigenvalues in $[0,1]$, so
$0\leq\rho'^2\leq\rho'$. For the positive operator
$K=(S-T)^\dagger(S-T)$, positivity of the trace pairing gives
$\Re\tr(K\rho'^2)\leq\Re\tr(K\rho')$. Thus the proved standard-distance
bound implies the literal-distance bound with the same coefficient $2$,
without changing either definition. The pure-state counterexamples below
also show that this coefficient cannot be decreased under the literal
reading. This comparison is a mathematical consequence of the proved
standard-distance theorem, not a separate claimed Lean declaration.

The Fourier--Naimark proof is consistent with the measurement normalization.
It proves
\[
  \mathbb E_a\operatorname{CON}_{\rho'}(A(a),\mathcal A(a))
  = \frac12+\frac12\mathbb E_{a,b}
       \tr_\rho(A(a)A(b)A(a+b)).
\]
The hypothesis therefore makes the consistency deficit at most $\delta/2$.
For binary projective measurements, the squared measurement distance is twice
that deficit, giving the bound $\delta$ for the squared binary-measurement
distance and the bound $2\delta$ for the squared state-dependent operator
distance.

\section{A four-point counterexample}

The bound $\delta$ for the mean squared state-dependent operator distance is
false.  Let $G=(\mathbb Z/2\mathbb Z)^2$, choose $r\in G\setminus\{0\}$, take
the one-dimensional state $\rho=1$, and define scalar binary observables by
\[
  A(a)=\begin{cases}
    -1,&a=r,\\
     1,&a\neq r.
  \end{cases}
\]
For the characters $\chi_u(a)=(-1)^{u\cdot a}$, the normalized Fourier
coefficients satisfy
\[
  \widehat A(0)=\frac12,
  \qquad
  \widehat A(u)=-\frac12\chi_u(r)\quad(u\neq0).
\]
Their multiset is
$\{\frac12,\frac12,\frac12,-\frac12\}$, and hence
\[
  \mathbb E_{a,b}A(a)A(b)A(a+b)
    =\sum_u\widehat A(u)^3=\frac14.
\]
The correlation hypothesis therefore holds with $\delta=3/4$.

Every scalar exact Boolean representation is one of the characters $\chi_u$,
and its mean squared state-dependent operator distance from $A$ is
\[
  \mathbb E_a |A(a)-\chi_u(a)|^2
    =2-2\widehat A(u)\geq1.
\]
Allowing an ancillary space does not reduce this lower bound.  Exact linearity
and the binary-observable relations make the operators $\mathcal A(a)$ a
commuting family of Hermitian involutions.  Their simultaneous spectral
subspaces carry characters of $G$.  Evaluation in any density operator is
therefore a convex combination of the scalar character correlations above, so
the mean squared state-dependent operator distance remains at least $1$.  It
cannot be at most $\delta=3/4$.  By~\eqref{eq:normalization}, the corresponding
mean squared binary-measurement distance has lower bound $1/2$, which is
compatible with the provider's conclusion.

There is also a direct positivity proof, independent of simultaneous
diagonalization. Choose $r=(1,1)$ and put $X=\mathcal A(1,0)$,
$Y=\mathcal A(0,1)$. Exact linearity gives $\mathcal A(0)=\Id$,
$\mathcal A(1,1)=XY=YX$, and
\[
 ((\Id-X)(\Id-Y))^\dagger((\Id-X)(\Id-Y))
   =4(\Id-X-Y+XY)\geq0.
\]
Thus the averaged correlation with $A$ is at most $1/2$ in every density
operator, giving operator error at least $1$. The constant identity
representation attains $1$, whereas the Fourier--Naimark construction has
error $2(1-C)=3/2$. The example therefore establishes an existence lower
bound of $4\delta/3$, not $2\delta$. These facts, and the negations of both
printed readings, are proved in \Lean{}.%
\footnote{See \leanid{blrCorrelation\_fourPointObservable},
\leanid{one\_le\_avg\_fourPointObservable\_distance},
\leanid{avg\_fourPointObservable\_distance\_one},
\leanid{not\_printedLinearityClaim}, and
\leanid{not\_printedLinearityLiteralDistanceClaim} in
\doclink{MIPStarRE/QPBT/Combining/Linearity/PrintedClaims.lean}.}

\section{Optimal universal constant}

For $k\geq1$, let $x,y\in\mathbb F_2^k$ and set
$f_k(x,y)=(-1)^{x\cdot y}$, acting as scalar observables on $\mathbb C$.
For $u,v\in\mathbb F_2^k$, character orthogonality gives
\[
 \widehat f_k(u,v)
 =\mathbb E_x(-1)^{u\cdot x}
       \mathbb E_y(-1)^{(x+v)\cdot y}
 =2^{-k}(-1)^{u\cdot v}.
\]
Every coefficient has magnitude $2^{-k}$, and their sum is $f_k(0,0)=1$.
The cubic identity in the provider proof therefore yields
$C_k=\sum_{u,v}\widehat f_k(u,v)^3=4^{-k}$.
Any exactly linear binary family on any finite ancillary space decomposes
into characters: equivalently, its Fourier coefficients are mutually
orthogonal projections summing to the identity. Their expectations in the
ancillary density operator are nonnegative weights summing to one. Its
correlation with $f_k$ is consequently at most
$\max_{u,v}\widehat f_k(u,v)=2^{-k}$. A character attaining this maximum
achieves equality, so the best existence error, even allowing all ancillas, is
\[
 E_{\mathrm{opt}}(f_k)=2(1-2^{-k}),\qquad
 \delta_k=1-4^{-k},\qquad
 \frac{E_{\mathrm{opt}}(f_k)}{\delta_k}
   =\frac{2}{1+2^{-k}}\longrightarrow2.
\]
Hence no constant $c<2$ can replace $2$ in a universal conclusion
$E_{\mathrm{opt}}\leq c\delta$ with the printed hypotheses. This proof
establishes minimality among constant-multiple repairs of the printed bound.
It does not determine the best bound as a function of each fixed $t$ and
$\delta$, or assert optimality of the particular Fourier-square construction.
The sharp normalization identity, that construction's exact error $2(1-C)$,
and the optimal existence error are three different assertions.

The argument covers an unbounded sequence of dimensions and is mathematical
evidence, not a claim of a formalized sharpness theorem. As additional finite
checks, exact rational enumeration of every scalar Boolean family for
$0\leq t\leq4$ found no violation of $E_{\mathrm{opt}}\leq2(1-C)$;
the quadratic examples for $1\leq k\leq5$ have respectively the ratios
$4/3,8/5,16/9,32/17,64/33$. The proof also covers mixed states and
noncommuting input observables through the Fourier--Naimark construction.
At zero deficit it gives zero error; for $t=0$ the same construction and
identities still apply. A negative $\delta$ cannot satisfy the hypothesis,
since the defect identity and positivity imply $C\leq1$.

\section{Comparison with the blueprint and formalization}

The QPBT paper quotes the quantum linearity theorem at lines~711--725 of its
Pauli-basis analysis and defers its proof to Natarajan and Vidick.%
\footnote{See
\path{references/qpbt-paper/14_analysis_of_the_pauli_basis_test.tex},
lines~711--725.}
The active blueprint expands the conclusion's distance as the state-dependent
operator distance $d$ and states the bound $2\delta$, together with the
equivalent bound $\delta$ for $\dbin$.  The formal theorem
\leanid{exists\_exactly\_linear\_observables} bounds the average over $u$
of \leanid{stateDepDistSq} by $2\delta$, and its companion
\leanid{exists\_exactly\_linear\_observables\_binaryObservableDistSq}
bounds the average of \leanid{binaryObservableDistSq} by $\delta$; the
printed bound $\delta$ for the mean squared state-dependent operator
distance is not reproduced.%
\footnote{See
\path{blueprint/src/chapter/ch15_qpbt_combining.tex},
\texttt{thm:linearity}, and
the two declarations in
\doclink{MIPStarRE/QPBT/Combining/Linearity.lean}.}

The mathematical content of the theorem is nevertheless available in the
correct normalization. The stability theorem assumes the mean squared
state-dependent multiplicative defect
\[
  \mathbb E_{a,b}\, d_\rho\big(A(a)A(b), A(a+b)\big)^2 \le \eta
\]
and produces an ancillary space, a unit vector $\ket{\mathrm{anc}}$ in it and
an exactly linear family $\{\mathcal A(a)\}$ of binary observables with
$\mathbb E_a\, d_{\rho'}(\mathcal A(a), A(a)\otimes\Id)^2 \le \eta$, where
$\rho' = \rho \otimes \ket{\mathrm{anc}}\bra{\mathrm{anc}}$.
Its bound rests on an identity rather than an inequality:
\leanid{avg\_stateDepDistSq\_roundedObservable\_eq\_avg\_multiplicativeDefect}
proves that the two averages above coincide, both being $2-2C$ for the
two-query correlation
$C = \mathbb E_{a,b}\Re\tr(A(a)A(b)A(a+b)\rho)$ of the linearity test.%
\footnote{The stability theorem is
\leanid{exists\_exact\_boolean\_representation}. Both declarations are in
\doclink{MIPStarRE/QPBT/Combining/Linearity/Stability.lean}; the blueprint
records them as \texttt{lem:boolean-representation-stability}.}
The source hypothesis $C \ge 1-\delta$ is therefore equivalent to the defect
bound with $\eta = 2\delta$. For positive $t$ and nonnegative $\delta$, the
stability theorem then yields the mean squared state-dependent operator
distance bound $2\delta$, without an upper bound on $\delta$ or a commutation
assumption on the $A(a)$. This reparameterization is the entire
content of the repair, and the corrected source-facing statement
\leanid{exists\_exactly\_linear\_observables} is obtained from
\leanid{exists\_exact\_boolean\_representation} by substituting
$\eta = 2\delta$.

The additional theorem
\leanid{exists\_exactly\_linear\_observables\_of\_correlation} proves this
same corrected conclusion for every $t\in\mathbb N$ and every real
$\delta$, without the unused numerical restrictions in the earlier API.
The universal propositions \leanid{PrintedLinearityClaim} and
\leanid{PrintedLinearityLiteralDistanceClaim} retain, respectively, the
single-density reading and the literal duplicated-density reading of the
printed theorem. Neither is asserted or supplied as a theorem hypothesis.
The second uses $d_{\rho'^2}$ in the standard notation; for the pure-state
four-point obstruction, $\rho'^2=\rho'$ on every allowed extension, so both
readings are refuted. This retention neither fixes a common ancilla in
advance nor restricts the printed existential domain.

The Boolean Fourier foundation keeps the state-dependent operator distance
unchanged and adds the separately named binary-measurement distance together
with equations~\eqref{eq:normalization}--\eqref{eq:binary-correlation}.  The
statement repair changed the conclusion to mean squared state-dependent
operator distance at most $2\delta$, added the equivalent form with mean
squared binary-measurement distance at most $\delta$, and synchronized the
blueprint; the realigned theorem is proved by the Fourier--Naimark argument,
with the distance/defect relation established as an equality, that is, by the
stability theorem under the reparameterization just described.%
\footnote{The normalization investigation was tracked in
\ghissue{124}; the source-facing statement repair was carried out in
\ghissue{129}.}

\section{Uses of the corrected estimate}

The provider applies Theorem~10 to the family $C(a,b)$ in its proof of
\texttt{thm:isotwoplayer}, lines~1537--1541. The resulting squared error is
$O(\varepsilon^{1/8})$; multiplying its constant by two preserves that
statement and all subsequent substitutions. The exact representation
relations are unchanged. The introduction to \texttt{thm:qblr\_game} at
lines~1115--1117 only promises the correlation hypothesis and does not use
the erroneous conclusion. The later EPR and energy applications consume
the same exact relations and unspecified robustness constants.

In QPBT, the sole invocation is chapter~14, line~832. For each fixed $(x,z)$
let $\eta_{x,z}$ be the conditional mean squared multiplicative defect.
The defect identity gives $C_{x,z}=1-\eta_{x,z}/2$, so applying the corrected
theorem with $\delta_{x,z}=\eta_{x,z}/2$ gives output error at most
$\eta_{x,z}$. Averaging again preserves the incoming bound. Exact linearity
and the value at zero give the projectors and completeness at lines~834--849;
Parseval and the averaged comparisons give self-consistency and both product
orders at lines~851--875 with the same polynomial error form.
These observations establish sufficiency of the normalization correction at
each use. They do not justify the separate assumption at line~832 that the
optional ancilla has already been absorbed into the fixed strategy.

The active blueprint's combined-point theorem has a direct field-valued
proof. No \verb|\uses{}| edge leaves the linearity theorem, import remark,
normalization remark, or common-ancilla lemma. The remaining references to
these nodes explain their proof or distinguish the direct point construction;
they are enumerated, together with the source applications, in the earlier
audit. In the current \Lean{} development, the operator-bound theorem is
called only by its binary-measurement companion. The unrestricted and
common-ancilla forms use the rounding identity directly and have no external
callers. The combined-point and extended-point constructions instead use the
independent field-valued measurement construction.%
\footnote{\path{audits/2026-09-22_linearity-normalization.md}, issue \#694.
The sole call to \leanid{exists\_exactly\_linear\_observables} is in
\leanid{exists\_exactly\_linear\_observables\_binaryObservableDistSq};
the unrestricted and common-ancilla forms are in the same
\doclink{MIPStarRE/QPBT/Combining/Linearity.lean}.
The independent declarations are \leanid{exists\_combinedPointsWitness} and
\leanid{exists\_extendedQ} in
\doclink{MIPStarRE/QPBT/Combining/Points.lean}.
The separate quotation and padding note is
\path{docs/paper-gaps/qpbt_linearity-theorem-quotation.tex}.}

\section{Conclusion}

The provider's binary-measurement calculation is correct through the factor
$1/2$ in equation~(3), but its final identification with the squared
state-dependent operator distance is incorrect.  The theorem proved by the
Fourier--Naimark argument has mean squared binary-measurement error at most
$\delta$, equivalently mean squared state-dependent operator error at most
$2\delta$.  The four-point example rules out the printed bound $\delta$ for
the mean squared state-dependent operator distance even when arbitrary
ancillary extensions are allowed.  The blueprint theorem and the formal
declarations carry the corrected bounds, including the unrestricted numerical
domain. The quadratic family proves optimality of the universal coefficient
$2$, but does not make the construction optimal for each input. The
normalization row is terminal as \texttt{documented-deviation}: the corrected
bounds are proved, while the printed operator bound remains refuted and its
two propositions remain unasserted. This documentary closure does not settle
state-preserving padding or certify other source-to-formalization comparisons.

\bibliographystyle{alpha}
\bibliography{references}

\end{document}
